% --- Archon physics formalization source begin ---
% archon:physics
% archon:covers PhyXMiniProblems/problem_phyx_mini_0335.lean
% archon:source-report reports/phyx_mini/problem_phyx_mini_0335.source.json

\chapter{Physics problem phyx\_mini\_0335}
\label{ch:PhyXMiniProblems_problem_phyx_mini_0335}

\paragraph{Problem source.}
\#\# Physical scenario

We can crudely model a nitrogen molecule as a pair of small balls, each with the mass of a nitrogen atom, \$2.3 \textbackslash{}times 10\textasciicircum{}\{-26\}\textbackslash{},\textbackslash{}text\{kg\}\$, attached by a rigid massless rod with length \$d = 2r = 188\textbackslash{},\textbackslash{}text\{pm\}\$.

\#\# Question

What is the moment of inertia of this molecule with respect to an axis passing through the midpoint and perpendicular to the molecular axis?

\#\# Answer choices

- A: A: 4.1 \textbackslash{}times 10\textasciicircum{}\{-46\}\textbackslash{}

- B: B: 5.1 \textbackslash{}times 10\textasciicircum{}\{-46\}\textbackslash{}

- C: C: 4.7 \textbackslash{}times 10\textasciicircum{}\{-46\}\textbackslash{}

- D: D: 6.1 \textbackslash{}times 10\textasciicircum{}\{-46\}\textbackslash{}

\#\# Figure caption (auxiliary; use the image as primary evidence)

The image consists of two diagrams labeled "BEFORE" and "AFTER," illustrating a physical concept involving motion.

**Before:**
- Two vertically aligned pairs of orange spheres are connected by gray rods.
- Green arrows labeled "v\_i" point horizontally to the right on the top pair and to the left on the bottom pair, indicating initial velocity.

**After:**
- Both pairs of orange spheres are shown at an angle, no longer vertical.
- Curved arrows indicate rotational motion around the center of the rods, labeled with the Greek letter "ω" (omega) to signify angular velocity.

Dashed lines separate each pair of spheres horizontally, suggesting a partition or axis of symmetry. The illustration likely represents a change from linear to rotational motion.

\#\# Dataset metadata

Ideal Gases and Kinetic Theory; Physical Model Grounding Reasoning, Multi-Formula Reasoning

\paragraph{Recorded answer/context.}
A: A: 4.1 \textbackslash{}times 10\textasciicircum{}\{-46\}\textbackslash{}

\paragraph{Figure/image path.}
/root/proposal\_for\_physic/science-mango/phyx\_mini\_run/phyx\_data/test\_image/335.png

\paragraph{Formalization target.}
create a compiling Lean file with sorry bodies at `PhyXMiniProblems/problem\_phyx\_mini\_0335.lean`. The Lean declarations must preserve the physical quantities, dimensions or dimensional roles, figure labels, governing-law hypotheses, and final relation expressed by this problem.
Use Mathlib/Physlib names found through LeanExplore where available. If a physics API is missing, introduce faithful local abstractions rather than scalar placeholder aliases.

\begin{theorem}[Physics formalization target]
\label{thm:physics:phyx_mini_0335:target}
\lean{PhyXMiniProblems.ProblemPhyXMini0335.nitrogen_molecule_moment_of_inertia}
\uses{def:physics:phyx-mini-0335:phyxminiproblems-problemphyxmini0335-momentofinertiainkilogrammeterssquared, def:physics:phyx-mini-0335:phyxminiproblems-problemphyxmini0335-nitrogenmoleculesetup, def:physics:phyx-mini-0335:phyxminiproblems-problemphyxmini0335-matchesnitrogenmoleculescenario, def:physics:phyx-mini-0335:phyxminiproblems-problemphyxmini0335-matchesproblemreadouts, def:physics:phyx-mini-0335:phyxminiproblems-problemphyxmini0335-matchesmidpointaxisgeometry, def:physics:phyx-mini-0335:phyxminiproblems-problemphyxmini0335-matchessuppliedbeforeafterfigure, def:physics:phyx-mini-0335:phyxminiproblems-problemphyxmini0335-hasphysicalnitrogenmoleculeparameters, def:physics:phyx-mini-0335:phyxminiproblems-problemphyxmini0335-satisfiestwopointmassinertialaw, def:physics:phyx-mini-0335:phyxminiproblems-problemphyxmini0335-answerchoice, def:physics:phyx-mini-0335:phyxminiproblems-problemphyxmini0335-displayedmomentofinertia, def:physics:phyx-mini-0335:phyxminiproblems-problemphyxmini0335-matchesanswerchoice}
The assigned autoformalize agent should translate this physics problem into Lean declarations in the covered file, with theorem and lemma proof bodies written as `by sorry`.
\end{theorem}
\begin{proof}
This is an autoformalization task, not a proof task. Produce statements that can later be proved by the physics prover without weakening the physical model.
\end{proof}

% --- Archon declaration topology begin ---
\section{Declaration topology}
\begin{definition}[speed Dimension]
  \label{def:physics:phyx-mini-0335:phyxminiproblems-problemphyxmini0335-speeddimension}
  \lean{PhyXMiniProblems.ProblemPhyXMini0335.speedDimension}
  The physical dimension of translational speed
\end{definition}
\begin{proof}
  This is the declared model definition of speed Dimension.
\end{proof}
\begin{definition}[angular Speed Dimension]
  \label{def:physics:phyx-mini-0335:phyxminiproblems-problemphyxmini0335-angularspeeddimension}
  \lean{PhyXMiniProblems.ProblemPhyXMini0335.angularSpeedDimension}
  The physical dimension of angular speed; radians are dimensionless
\end{definition}
\begin{proof}
  This is the declared model definition of angular Speed Dimension.
\end{proof}
\begin{definition}[moment Of Inertia Dimension]
  \label{def:physics:phyx-mini-0335:phyxminiproblems-problemphyxmini0335-momentofinertiadimension}
  \lean{PhyXMiniProblems.ProblemPhyXMini0335.momentOfInertiaDimension}
  The physical dimension of moment of inertia, mass times length squared
\end{definition}
\begin{proof}
  This is the declared model definition of moment Of Inertia Dimension.
\end{proof}
\begin{definition}[Mass Quantity]
  \label{def:physics:phyx-mini-0335:phyxminiproblems-problemphyxmini0335-massquantity}
  \lean{PhyXMiniProblems.ProblemPhyXMini0335.MassQuantity}
  A nonnegative physical mass, independent of its readout unit
\end{definition}
\begin{proof}
  This is the declared model definition of Mass Quantity.
\end{proof}
\begin{definition}[Length Quantity]
  \label{def:physics:phyx-mini-0335:phyxminiproblems-problemphyxmini0335-lengthquantity}
  \lean{PhyXMiniProblems.ProblemPhyXMini0335.LengthQuantity}
  A nonnegative physical length, independent of its readout unit
\end{definition}
\begin{proof}
  This is the declared model definition of Length Quantity.
\end{proof}
\begin{definition}[Speed Magnitude Quantity]
  \label{def:physics:phyx-mini-0335:phyxminiproblems-problemphyxmini0335-speedmagnitudequantity}
  \lean{PhyXMiniProblems.ProblemPhyXMini0335.SpeedMagnitudeQuantity}
  \uses{def:physics:phyx-mini-0335:phyxminiproblems-problemphyxmini0335-speeddimension}
  A nonnegative translational speed magnitude
\end{definition}
\begin{proof}
  This is the declared model definition of Speed Magnitude Quantity.
\end{proof}
\begin{definition}[Angular Speed Magnitude Quantity]
  \label{def:physics:phyx-mini-0335:phyxminiproblems-problemphyxmini0335-angularspeedmagnitudequantity}
  \lean{PhyXMiniProblems.ProblemPhyXMini0335.AngularSpeedMagnitudeQuantity}
  \uses{def:physics:phyx-mini-0335:phyxminiproblems-problemphyxmini0335-angularspeeddimension}
  A nonnegative angular-speed magnitude
\end{definition}
\begin{proof}
  This is the declared model definition of Angular Speed Magnitude Quantity.
\end{proof}
\begin{definition}[Moment Of Inertia Quantity]
  \label{def:physics:phyx-mini-0335:phyxminiproblems-problemphyxmini0335-momentofinertiaquantity}
  \lean{PhyXMiniProblems.ProblemPhyXMini0335.MomentOfInertiaQuantity}
  \uses{def:physics:phyx-mini-0335:phyxminiproblems-problemphyxmini0335-momentofinertiadimension}
  A nonnegative physical moment of inertia about a specified axis
\end{definition}
\begin{proof}
  This is the declared model definition of Moment Of Inertia Quantity.
\end{proof}
\begin{definition}[mass Readout]
  \label{def:physics:phyx-mini-0335:phyxminiproblems-problemphyxmini0335-massreadout}
  \lean{PhyXMiniProblems.ProblemPhyXMini0335.massReadout}
  \uses{def:physics:phyx-mini-0335:phyxminiproblems-problemphyxmini0335-massquantity}
  Read a physical mass in a selected mass unit
\end{definition}
\begin{proof}
  This is the declared model definition of mass Readout.
\end{proof}
\begin{definition}[length Readout]
  \label{def:physics:phyx-mini-0335:phyxminiproblems-problemphyxmini0335-lengthreadout}
  \lean{PhyXMiniProblems.ProblemPhyXMini0335.lengthReadout}
  \uses{def:physics:phyx-mini-0335:phyxminiproblems-problemphyxmini0335-lengthquantity}
  Read a physical length in a selected length unit
\end{definition}
\begin{proof}
  This is the declared model definition of length Readout.
\end{proof}
\begin{definition}[speed Readout]
  \label{def:physics:phyx-mini-0335:phyxminiproblems-problemphyxmini0335-speedreadout}
  \lean{PhyXMiniProblems.ProblemPhyXMini0335.speedReadout}
  \uses{def:physics:phyx-mini-0335:phyxminiproblems-problemphyxmini0335-speedmagnitudequantity}
  Read speed in coherent selected length and time units
\end{definition}
\begin{proof}
  This is the declared model definition of speed Readout.
\end{proof}
\begin{definition}[angular Speed Readout]
  \label{def:physics:phyx-mini-0335:phyxminiproblems-problemphyxmini0335-angularspeedreadout}
  \lean{PhyXMiniProblems.ProblemPhyXMini0335.angularSpeedReadout}
  \uses{def:physics:phyx-mini-0335:phyxminiproblems-problemphyxmini0335-angularspeedmagnitudequantity}
  Read angular speed in radians per selected time unit
\end{definition}
\begin{proof}
  This is the declared model definition of angular Speed Readout.
\end{proof}
\begin{definition}[moment Of Inertia Readout]
  \label{def:physics:phyx-mini-0335:phyxminiproblems-problemphyxmini0335-momentofinertiareadout}
  \lean{PhyXMiniProblems.ProblemPhyXMini0335.momentOfInertiaReadout}
  \uses{def:physics:phyx-mini-0335:phyxminiproblems-problemphyxmini0335-momentofinertiaquantity}
  Read moment of inertia in coherent selected mass and length units
\end{definition}
\begin{proof}
  This is the declared model definition of moment Of Inertia Readout.
\end{proof}
\begin{definition}[mass In Kilograms]
  \label{def:physics:phyx-mini-0335:phyxminiproblems-problemphyxmini0335-massinkilograms}
  \lean{PhyXMiniProblems.ProblemPhyXMini0335.massInKilograms}
  \uses{def:physics:phyx-mini-0335:phyxminiproblems-problemphyxmini0335-massquantity, def:physics:phyx-mini-0335:phyxminiproblems-problemphyxmini0335-massreadout}
  Kilogram readout of an atom's physical mass
\end{definition}
\begin{proof}
  This is the declared model definition of mass In Kilograms.
\end{proof}
\begin{definition}[length In Meters]
  \label{def:physics:phyx-mini-0335:phyxminiproblems-problemphyxmini0335-lengthinmeters}
  \lean{PhyXMiniProblems.ProblemPhyXMini0335.lengthInMeters}
  \uses{def:physics:phyx-mini-0335:phyxminiproblems-problemphyxmini0335-lengthquantity, def:physics:phyx-mini-0335:phyxminiproblems-problemphyxmini0335-lengthreadout}
  Meter readout used in the coherent SI inertia law
\end{definition}
\begin{proof}
  This is the declared model definition of length In Meters.
\end{proof}
\begin{definition}[length In Picometers]
  \label{def:physics:phyx-mini-0335:phyxminiproblems-problemphyxmini0335-lengthinpicometers}
  \lean{PhyXMiniProblems.ProblemPhyXMini0335.lengthInPicometers}
  \uses{def:physics:phyx-mini-0335:phyxminiproblems-problemphyxmini0335-lengthquantity, def:physics:phyx-mini-0335:phyxminiproblems-problemphyxmini0335-lengthreadout}
  Picometer readout used for the stated molecular separation
\end{definition}
\begin{proof}
  This is the declared model definition of length In Picometers.
\end{proof}
\begin{definition}[moment Of Inertia In Kilogram Meters Squared]
  \label{def:physics:phyx-mini-0335:phyxminiproblems-problemphyxmini0335-momentofinertiainkilogrammeterssquared}
  \lean{PhyXMiniProblems.ProblemPhyXMini0335.momentOfInertiaInKilogramMetersSquared}
  \uses{def:physics:phyx-mini-0335:phyxminiproblems-problemphyxmini0335-momentofinertiaquantity, def:physics:phyx-mini-0335:phyxminiproblems-problemphyxmini0335-momentofinertiareadout}
  SI readout of a moment of inertia, in kilogram-meter squared
\end{definition}
\begin{proof}
  This is the declared model definition of moment Of Inertia In Kilogram Meters Squared.
\end{proof}
\begin{definition}[Atom Site]
  \label{def:physics:phyx-mini-0335:phyxminiproblems-problemphyxmini0335-atomsite}
  \lean{PhyXMiniProblems.ProblemPhyXMini0335.AtomSite}
  The two nitrogen-atom sites at the ends of the molecular bond
\end{definition}
\begin{proof}
  This is the declared model definition of Atom Site.
\end{proof}
\begin{definition}[Molecular Species]
  \label{def:physics:phyx-mini-0335:phyxminiproblems-problemphyxmini0335-molecularspecies}
  \lean{PhyXMiniProblems.ProblemPhyXMini0335.MolecularSpecies}
  The molecular species stated in the problem
\end{definition}
\begin{proof}
  This is the declared model definition of Molecular Species.
\end{proof}
\begin{definition}[Atom Body Model]
  \label{def:physics:phyx-mini-0335:phyxminiproblems-problemphyxmini0335-atombodymodel}
  \lean{PhyXMiniProblems.ProblemPhyXMini0335.AtomBodyModel}
  The idealization of each orange ball as a small point mass
\end{definition}
\begin{proof}
  This is the declared model definition of Atom Body Model.
\end{proof}
\begin{definition}[Bond Rod Model]
  \label{def:physics:phyx-mini-0335:phyxminiproblems-problemphyxmini0335-bondrodmodel}
  \lean{PhyXMiniProblems.ProblemPhyXMini0335.BondRodModel}
  The idealization of the connector between the two atom sites
\end{definition}
\begin{proof}
  This is the declared model definition of Bond Rod Model.
\end{proof}
\begin{definition}[Rotation Axis Geometry]
  \label{def:physics:phyx-mini-0335:phyxminiproblems-problemphyxmini0335-rotationaxisgeometry}
  \lean{PhyXMiniProblems.ProblemPhyXMini0335.RotationAxisGeometry}
  Geometry of the axis about which the moment of inertia is requested
\end{definition}
\begin{proof}
  This is the declared model definition of Rotation Axis Geometry.
\end{proof}
\begin{definition}[Figure Panel]
  \label{def:physics:phyx-mini-0335:phyxminiproblems-problemphyxmini0335-figurepanel}
  \lean{PhyXMiniProblems.ProblemPhyXMini0335.FigurePanel}
  The two panels explicitly captioned in the supplied image
\end{definition}
\begin{proof}
  This is the declared model definition of Figure Panel.
\end{proof}
\begin{definition}[Figure Molecule]
  \label{def:physics:phyx-mini-0335:phyxminiproblems-problemphyxmini0335-figuremolecule}
  \lean{PhyXMiniProblems.ProblemPhyXMini0335.FigureMolecule}
  The upper and lower dumbbells separated by the dashed horizontal line
\end{definition}
\begin{proof}
  This is the declared model definition of Figure Molecule.
\end{proof}
\begin{definition}[Rod Orientation]
  \label{def:physics:phyx-mini-0335:phyxminiproblems-problemphyxmini0335-rodorientation}
  \lean{PhyXMiniProblems.ProblemPhyXMini0335.RodOrientation}
  Orientations of a dumbbell rod appearing in the raster image
\end{definition}
\begin{proof}
  This is the declared model definition of Rod Orientation.
\end{proof}
\begin{definition}[Horizontal Direction]
  \label{def:physics:phyx-mini-0335:phyxminiproblems-problemphyxmini0335-horizontaldirection}
  \lean{PhyXMiniProblems.ProblemPhyXMini0335.HorizontalDirection}
  Horizontal senses of the green arrows labelled v\_i
\end{definition}
\begin{proof}
  This is the declared model definition of Horizontal Direction.
\end{proof}
\begin{definition}[Rotation Sense]
  \label{def:physics:phyx-mini-0335:phyxminiproblems-problemphyxmini0335-rotationsense}
  \lean{PhyXMiniProblems.ProblemPhyXMini0335.RotationSense}
  Rotation senses of the curved arrows labelled omega
\end{definition}
\begin{proof}
  This is the declared model definition of Rotation Sense.
\end{proof}
\begin{definition}[Supplied Before After Figure]
  \label{def:physics:phyx-mini-0335:phyxminiproblems-problemphyxmini0335-suppliedbeforeafterfigure}
  \lean{PhyXMiniProblems.ProblemPhyXMini0335.SuppliedBeforeAfterFigure}
  \uses{def:physics:phyx-mini-0335:phyxminiproblems-problemphyxmini0335-speedmagnitudequantity, def:physics:phyx-mini-0335:phyxminiproblems-problemphyxmini0335-angularspeedmagnitudequantity, def:physics:phyx-mini-0335:phyxminiproblems-problemphyxmini0335-figurepanel, def:physics:phyx-mini-0335:phyxminiproblems-problemphyxmini0335-figuremolecule, def:physics:phyx-mini-0335:phyxminiproblems-problemphyxmini0335-rodorientation, def:physics:phyx-mini-0335:phyxminiproblems-problemphyxmini0335-horizontaldirection, def:physics:phyx-mini-0335:phyxminiproblems-problemphyxmini0335-rotationsense}
  Data represented in the supplied raster image. The image depicts two dumbbells before and after an interaction. The common v\_i and omega labels are retained as dimensionful speed magnitudes, even though neither is needed to calculate the molecule's static moment of inertia
\end{definition}
\begin{proof}
  This is the declared model definition of Supplied Before After Figure.
\end{proof}
\begin{definition}[Nitrogen Molecule Setup]
  \label{def:physics:phyx-mini-0335:phyxminiproblems-problemphyxmini0335-nitrogenmoleculesetup}
  \lean{PhyXMiniProblems.ProblemPhyXMini0335.NitrogenMoleculeSetup}
  \uses{def:physics:phyx-mini-0335:phyxminiproblems-problemphyxmini0335-massquantity, def:physics:phyx-mini-0335:phyxminiproblems-problemphyxmini0335-lengthquantity, def:physics:phyx-mini-0335:phyxminiproblems-problemphyxmini0335-momentofinertiaquantity, def:physics:phyx-mini-0335:phyxminiproblems-problemphyxmini0335-atomsite, def:physics:phyx-mini-0335:phyxminiproblems-problemphyxmini0335-molecularspecies, def:physics:phyx-mini-0335:phyxminiproblems-problemphyxmini0335-atombodymodel, def:physics:phyx-mini-0335:phyxminiproblems-problemphyxmini0335-bondrodmodel, def:physics:phyx-mini-0335:phyxminiproblems-problemphyxmini0335-rotationaxisgeometry, def:physics:phyx-mini-0335:phyxminiproblems-problemphyxmini0335-suppliedbeforeafterfigure}
  Independent physical quantities in the molecular model. In particular, the moment of inertia is an observable constrained only by the governing law below; it is not defined to equal a displayed answer
\end{definition}
\begin{proof}
  This is the declared model definition of Nitrogen Molecule Setup.
\end{proof}
\begin{definition}[Matches Nitrogen Molecule Scenario]
  \label{def:physics:phyx-mini-0335:phyxminiproblems-problemphyxmini0335-matchesnitrogenmoleculescenario}
  \lean{PhyXMiniProblems.ProblemPhyXMini0335.MatchesNitrogenMoleculeScenario}
  \uses{def:physics:phyx-mini-0335:phyxminiproblems-problemphyxmini0335-nitrogenmoleculesetup}
  The qualitative molecular and rotation-axis model stated in the prose
\end{definition}
\begin{proof}
  This is the declared model definition of Matches Nitrogen Molecule Scenario.
\end{proof}
\begin{definition}[Matches Problem Readouts]
  \label{def:physics:phyx-mini-0335:phyxminiproblems-problemphyxmini0335-matchesproblemreadouts}
  \lean{PhyXMiniProblems.ProblemPhyXMini0335.MatchesProblemReadouts}
  \uses{def:physics:phyx-mini-0335:phyxminiproblems-problemphyxmini0335-massinkilograms, def:physics:phyx-mini-0335:phyxminiproblems-problemphyxmini0335-lengthinpicometers, def:physics:phyx-mini-0335:phyxminiproblems-problemphyxmini0335-atomsite, def:physics:phyx-mini-0335:phyxminiproblems-problemphyxmini0335-nitrogenmoleculesetup}
  The numerical data printed in the question: each atom has mass 2.3 * 10\textasciicircum{}-26 kg, and the atom-to-atom separation is 188 pm. The requested moment of inertia does not occur in these readouts
\end{definition}
\begin{proof}
  This is the declared model definition of Matches Problem Readouts.
\end{proof}
\begin{definition}[Matches Midpoint Axis Geometry]
  \label{def:physics:phyx-mini-0335:phyxminiproblems-problemphyxmini0335-matchesmidpointaxisgeometry}
  \lean{PhyXMiniProblems.ProblemPhyXMini0335.MatchesMidpointAxisGeometry}
  \uses{def:physics:phyx-mini-0335:phyxminiproblems-problemphyxmini0335-lengthreadout, def:physics:phyx-mini-0335:phyxminiproblems-problemphyxmini0335-atomsite, def:physics:phyx-mini-0335:phyxminiproblems-problemphyxmini0335-nitrogenmoleculesetup}
  The stated relation d = 2 r and the midpoint-axis geometry. Both atom sites are one half-separation from the perpendicular axis. These are geometric inputs, not the requested inertia value
\end{definition}
\begin{proof}
  This is the declared model definition of Matches Midpoint Axis Geometry.
\end{proof}
\begin{definition}[Matches Supplied Before After Figure]
  \label{def:physics:phyx-mini-0335:phyxminiproblems-problemphyxmini0335-matchessuppliedbeforeafterfigure}
  \lean{PhyXMiniProblems.ProblemPhyXMini0335.MatchesSuppliedBeforeAfterFigure}
  \uses{def:physics:phyx-mini-0335:phyxminiproblems-problemphyxmini0335-figurepanel, def:physics:phyx-mini-0335:phyxminiproblems-problemphyxmini0335-figuremolecule, def:physics:phyx-mini-0335:phyxminiproblems-problemphyxmini0335-nitrogenmoleculesetup}
  Literal panel captions, arrow labels, orientations, and motion senses read from the primary image. The upper initial arrow points right, the lower one points left, and the two after-panel angular arrows have opposite senses. The dashed divider is recorded without interpreting it as the rotation axis
\end{definition}
\begin{proof}
  This is the declared model definition of Matches Supplied Before After Figure.
\end{proof}
\begin{definition}[Has Physical Nitrogen Molecule Parameters]
  \label{def:physics:phyx-mini-0335:phyxminiproblems-problemphyxmini0335-hasphysicalnitrogenmoleculeparameters}
  \lean{PhyXMiniProblems.ProblemPhyXMini0335.HasPhysicalNitrogenMoleculeParameters}
  \uses{def:physics:phyx-mini-0335:phyxminiproblems-problemphyxmini0335-speedreadout, def:physics:phyx-mini-0335:phyxminiproblems-problemphyxmini0335-angularspeedreadout, def:physics:phyx-mini-0335:phyxminiproblems-problemphyxmini0335-massinkilograms, def:physics:phyx-mini-0335:phyxminiproblems-problemphyxmini0335-lengthinmeters, def:physics:phyx-mini-0335:phyxminiproblems-problemphyxmini0335-atomsite, def:physics:phyx-mini-0335:phyxminiproblems-problemphyxmini0335-nitrogenmoleculesetup}
  Positivity and nondegeneracy conditions for the modeled quantities
\end{definition}
\begin{proof}
  This is the declared model definition of Has Physical Nitrogen Molecule Parameters.
\end{proof}
\begin{definition}[Satisfies Two Point Mass Inertia Law]
  \label{def:physics:phyx-mini-0335:phyxminiproblems-problemphyxmini0335-satisfiestwopointmassinertialaw}
  \lean{PhyXMiniProblems.ProblemPhyXMini0335.SatisfiesTwoPointMassInertiaLaw}
  \uses{def:physics:phyx-mini-0335:phyxminiproblems-problemphyxmini0335-massreadout, def:physics:phyx-mini-0335:phyxminiproblems-problemphyxmini0335-lengthreadout, def:physics:phyx-mini-0335:phyxminiproblems-problemphyxmini0335-momentofinertiareadout, def:physics:phyx-mini-0335:phyxminiproblems-problemphyxmini0335-atomsite, def:physics:phyx-mini-0335:phyxminiproblems-problemphyxmini0335-nitrogenmoleculesetup}
  The general point-mass law I = sum m\_i rho\_i\textasciicircum{}2, stated in every coherent mass and length unit. Its radii are independent setup observables related to d / 2 only by the separate geometry hypotheses. Thus this law contains neither the simplified molecular formula nor the numerical answer
\end{definition}
\begin{proof}
  This is the declared model definition of Satisfies Two Point Mass Inertia Law.
\end{proof}
\begin{definition}[Answer Choice]
  \label{def:physics:phyx-mini-0335:phyxminiproblems-problemphyxmini0335-answerchoice}
  \lean{PhyXMiniProblems.ProblemPhyXMini0335.AnswerChoice}
  The four answer labels printed beside the problem
\end{definition}
\begin{proof}
  This is the declared model definition of Answer Choice.
\end{proof}
\begin{definition}[displayed Moment Of Inertia]
  \label{def:physics:phyx-mini-0335:phyxminiproblems-problemphyxmini0335-displayedmomentofinertia}
  \lean{PhyXMiniProblems.ProblemPhyXMini0335.displayedMomentOfInertia}
  \uses{def:physics:phyx-mini-0335:phyxminiproblems-problemphyxmini0335-answerchoice}
  Displayed moments of inertia, read in kilogram-meter squared
\end{definition}
\begin{proof}
  This is the declared model definition of displayed Moment Of Inertia.
\end{proof}
\begin{definition}[displayed Resolution]
  \label{def:physics:phyx-mini-0335:phyxminiproblems-problemphyxmini0335-displayedresolution}
  \lean{PhyXMiniProblems.ProblemPhyXMini0335.displayedResolution}
  One unit in the final displayed digit of every answer choice
\end{definition}
\begin{proof}
  This is the declared model definition of displayed Resolution.
\end{proof}
\begin{definition}[Matches Answer Choice]
  \label{def:physics:phyx-mini-0335:phyxminiproblems-problemphyxmini0335-matchesanswerchoice}
  \lean{PhyXMiniProblems.ProblemPhyXMini0335.MatchesAnswerChoice}
  \uses{def:physics:phyx-mini-0335:phyxminiproblems-problemphyxmini0335-momentofinertiainkilogrammeterssquared, def:physics:phyx-mini-0335:phyxminiproblems-problemphyxmini0335-nitrogenmoleculesetup, def:physics:phyx-mini-0335:phyxminiproblems-problemphyxmini0335-answerchoice, def:physics:phyx-mini-0335:phyxminiproblems-problemphyxmini0335-displayedmomentofinertia, def:physics:phyx-mini-0335:phyxminiproblems-problemphyxmini0335-displayedresolution}
  Rounding agreement with a displayed answer choice
\end{definition}
\begin{proof}
  This is the declared model definition of Matches Answer Choice.
\end{proof}
\begin{lemma}[midpoint radius in meters]
  \label{lem:physics:phyx-mini-0335:phyxminiproblems-problemphyxmini0335-midpoint-radius-in-meters}
  \lean{PhyXMiniProblems.ProblemPhyXMini0335.midpoint_radius_in_meters}
  \uses{def:physics:phyx-mini-0335:phyxminiproblems-problemphyxmini0335-lengthinmeters, def:physics:phyx-mini-0335:phyxminiproblems-problemphyxmini0335-nitrogenmoleculesetup, def:physics:phyx-mini-0335:phyxminiproblems-problemphyxmini0335-matchesmidpointaxisgeometry}
  The midpoint relation gives r = d / 2 in meter readouts
\end{lemma}
\begin{proof}
  The conclusion follows from the governing relations and figure readouts named in the statement of midpoint radius in meters.
\end{proof}
\begin{lemma}[diatomic point mass inertia formula]
  \label{lem:physics:phyx-mini-0335:phyxminiproblems-problemphyxmini0335-diatomic-point-mass-inertia-formula}
  \lean{PhyXMiniProblems.ProblemPhyXMini0335.diatomic_point_mass_inertia_formula}
  \uses{def:physics:phyx-mini-0335:phyxminiproblems-problemphyxmini0335-massinkilograms, def:physics:phyx-mini-0335:phyxminiproblems-problemphyxmini0335-lengthinmeters, def:physics:phyx-mini-0335:phyxminiproblems-problemphyxmini0335-momentofinertiainkilogrammeterssquared, def:physics:phyx-mini-0335:phyxminiproblems-problemphyxmini0335-nitrogenmoleculesetup, def:physics:phyx-mini-0335:phyxminiproblems-problemphyxmini0335-matchesproblemreadouts, def:physics:phyx-mini-0335:phyxminiproblems-problemphyxmini0335-matchesmidpointaxisgeometry, def:physics:phyx-mini-0335:phyxminiproblems-problemphyxmini0335-satisfiestwopointmassinertialaw}
  The two equal point masses and the midpoint geometry give the familiar diatomic formula I = m d\textasciicircum{}2 / 2 in SI readouts. This is derived rather than assumed by the governing-law interface
\end{lemma}
\begin{proof}
  The conclusion follows from the governing relations and figure readouts named in the statement of diatomic point mass inertia formula.
\end{proof}
% --- Archon declaration topology end ---
% --- Archon physics formalization source end ---
